\subsection{Ablations}
The main results suggest that a macro commit helps only when it satisfies the three conditions built into the commit rule of Section~\ref{sec:commit}. 
The reused steps must stay hidden from the model, since asking the model to select mined patterns adds a decision problem rather than skipping any actor call; the runtime must filter mined candidates aggressively, since a library match alone weakly predicts that the trajectory will follow the mined pattern; and a commit must skip enough actor calls on the critical path to overcome speculation, verification, and synchronization overhead. 
The three ablations below test these conditions in turn.

\subsection{Interface: hidden runtime state}
 
Table~\ref{tab:single-gpu-alternatives} evaluates two simpler ways of using mined routines on a single GPU. 
The first exposes mined patterns as registered meta-tools~\citep{awo-metatools-2026}, making reuse a model-visible action. 
The second leaves the model interface unchanged but commits matched patterns passively, without the online checks of the full commit rule.

Neither alternative has desirable results. 
Registered meta-tools slightly increase latency and lower accuracy, because the actor almost never selects them: exposing mined routines as ordinary tools does not reliably skip actor calls, it poses an extra tool-selection problem. 
Passive committing has the opposite failure mode. It cuts latency by 11.34\%, confirming that hidden runtime reuse can skip actor calls, but accuracy falls from 99.52\% to 96.48\%. 
Hidden reuse alone is therefore insufficient: the runtime must also verify that the committed steps are anchored to the authoritative trajectory. SMC keeps the latency benefit of hidden reuse while adding the online commit conditions.

\begin{table}[t]
\centering
\small
\caption{Results on $\tau^2$-Bench with single GPU setting. }
\label{tab:single-gpu-alternatives}
\begin{tabular}{lrrr}
\toprule
Run & Acc. & Lat. & $\Delta$(\%) \\
\midrule
Baseline & 99.52 & 27.60 & -- \\
Baseline + AWO-like & 99.34 & 27.89 & +1.05 \\
Baseline + passive & 96.48 & 24.47 & -11.34 \\
\bottomrule
\end{tabular}
\end{table}

\subsection{Commit Precision}

The passive-committing result raises the central safety question: how does SMC avoid committing incorrect mined patterns? Table~\ref{tab:filter-precision} answers with a staged audit on 100 held-out $\tau^2$ Telecom tasks.
% not used for mining. 
For the first four rows, SMC suppresses the commit and lets the actor continue, so the realized trajectory provides an event-level counterfactual for whether the candidate steps were exact; the final row reports task-outcome preservation for the steps SMC actually committed.

The mined library by itself is far too noisy to act as a commit rule: a library match reproduces the exact future steps in only 34.6\% of events. 
Requiring the drafter to have already executed the steps raises precision to 70.6\%, since SMC no longer commits patterns that exist only as offline matches. 
Verifying the anchor call raises it to 87.9\% by tying the candidate to the authoritative trajectory, and the depth guard ($L_{\min}=1$, at least one committed step after the anchor) raises it to 90.4\% while removing shallow commits unlikely to repay synchronization overhead. 
Under full rule, all committed events preserve the outcome in this audit. 
As noted in Section~\ref{sec:commit}, this is not a proof that every committed pattern is uniquely correct; it shows that the online checks turn a low-precision library into a high-precision set of commits on this held-out audit.

\begin{table}[t]
\centering
\small
\caption{Conditional precision of the commit stages on 100 held-out $\tau^2$ Telecom tasks. Rows above the rule are event-level exact-match checks; the committed row is a task-outcome check.}
\label{tab:filter-precision}
\begin{tabular}{lrr}
\toprule
Filter stage & Events & \makecell{Correct / \\outcome-preserving} \\
\midrule
Library match only & 1,968 & 681/1,968 (34.6\%) \\
+ drafter-executed & 885 & 625/885 (70.6\%) \\
+ verified anchor call & 711 & 625/711 (87.9\%) \\
+ depth guard ($L_{\min}=1$) & 343 & 310/343 (90.4\%) \\
\midrule
Actually fired & 158 & 158/158 (100.0\%) \\
\bottomrule
\end{tabular}
\end{table}

\subsection{Critical-Path Depth}

High precision alone does not guarantee speedup: a correct macro is unhelpful if it skips only one or two actor calls or skips calls that are off the critical path.
Table~\ref{tab:legacy-final} compares the final SMC runtime with an earlier runtime that committed more aggressively.
The legacy runtime commits almost twice as often as the final runtime and skips more total steps, yet it is 1.64\% slower than SA. 
Raw macro count is thus the wrong objective: many additional legacy commits are too shallow or poorly aligned with already executed draft work, adding commit overhead without reducing end-to-end latency. 
The final runtime suppresses 6,281 depth-one opportunities, enforces the $L_{\min}$ depth floor, and commits only steps already executed on the draft branch. 
While it commits fewer macros and skips fewer total steps, it is 10.23\% faster than SA.
The useful unit is therefore not a macro hit but a guarded, sufficiently deep commit that skips actor calls on the critical path.

\begin{table}[t]
\centering
\small
\caption{$\tau^2$ Telecom full-run comparison: raw macro count versus critical-path commits. Hits is the number of macro commits; Skip is the total committed draft steps after anchors.}
\label{tab:legacy-final}
\begin{tabular}{lrrrrr}
\toprule
Run & Acc. & Lat. & Hits & Skip & $\Delta$(\%) \\
\midrule
SA & 99.47 & 25.03 & 0 & 0 & -- \\
Legacy macro & 99.56 & 25.44 & 6,410 & 10,528 & +1.64 \\
Final SMC & 99.52 & 22.47 & 3,352 & 7,154 & -10.23 \\
\bottomrule
\end{tabular}
\end{table}